% =====================================================================
% main.tex
% Dokumen induk Buku Ajar Teknik Kompilasi.
% Bab dimuat dengan \subfile agar setiap chapter.tex tetap dapat
% dikompilasi mandiri.
% =====================================================================

\documentclass[12pt,a4paper,oneside]{book}

% =====================================================================
% preamble.tex
% Urutan pemuatan konfigurasi bersifat tetap.
% =====================================================================

% =====================================================================
% config/packages.tex
% Semua paket LaTeX dimuat di sini.
% Bab TIDAK BOLEH memuat \usepackage (Aturan 3 Arsitektur Buku Ajar).
% =====================================================================

% --- Encoding dan bahasa ---------------------------------------------
\usepackage[utf8]{inputenc}
\usepackage[T1]{fontenc}
\usepackage[indonesian]{babel}
\usepackage{lmodern}
\usepackage{textcomp}
\usepackage{microtype}

% --- Tata letak halaman ----------------------------------------------
\usepackage[
  a4paper,
  left=2.8cm,
  right=2.5cm,
  top=2.8cm,
  bottom=3.0cm,
  headheight=15pt,
  headsep=0.5cm,
  footskip=1.0cm
]{geometry}
\usepackage{fancyhdr}
\usepackage{emptypage}

% --- Matematika -------------------------------------------------------
\usepackage{amsmath}
\usepackage{amssymb}
\usepackage{amsthm}
\usepackage{mathtools}

% --- Grafik dan diagram ----------------------------------------------
\usepackage{graphicx}
\usepackage{float}
\usepackage{tikz}
\usepackage{forest}

% --- Tabel ------------------------------------------------------------
\usepackage{booktabs}
\usepackage{longtable}
\usepackage{array}
\usepackage{tabularx}
\usepackage{multirow}
% colortbl dimuat hanya bila benar-benar diperlukan; lihat catatan di
% bawah. Paket ini menambal internal array/tabularx dan dapat memicu
% "Undefined control sequence \@nextchar" pada tabularx.
% \usepackage{colortbl}

% --- Daftar -----------------------------------------------------------
\usepackage{enumitem}

% --- Warna dan kotak --------------------------------------------------
\usepackage{xcolor}
\usepackage[most]{tcolorbox}

% --- Kode program -----------------------------------------------------
\usepackage{listings}
\usepackage{upquote}

% --- Keterangan dan navigasi -----------------------------------------
\usepackage{caption}
\usepackage{tocloft}
\usepackage{titlesec}
\usepackage{placeins}

% --- Referensi dan tautan --------------------------------------------
\usepackage[numbers,sort&compress]{natbib}
\usepackage{url}
\usepackage{xurl}
\usepackage[unicode,breaklinks=true,pdfencoding=auto]{hyperref}

% --- Utilitas ---------------------------------------------------------
\usepackage{etoolbox}
\usepackage{needspace}

% --- Kompilasi modular (dimuat terakhir) -----------------------------
\usepackage{subfiles}

% =====================================================================
% config/styles.tex
% Warna, tata letak, formato judul, header/footer, dan penomoran.
% =====================================================================

% --- Palet warna ------------------------------------------------------
\definecolor{ubbblue}{RGB}{15,55,110}
\definecolor{ubbteal}{RGB}{0,120,130}
\definecolor{ubbgold}{RGB}{200,150,40}
\definecolor{codebg}{RGB}{247,247,247}
\definecolor{codeframe}{RGB}{210,215,225}
\definecolor{textgray}{RGB}{90,90,90}

% --- Spasi dan paragraf ------------------------------------------------
\setlength{\parindent}{0pt}
\setlength{\parskip}{0.5em}
\linespread{1.12}
\setlength{\emergencystretch}{2em}
\raggedbottom
\clubpenalty=1000
\widowpenalty=1000
\renewcommand{\arraystretch}{1.15}

% --- Format judul bab --------------------------------------------------
\titleformat{\chapter}[display]
  {\normalfont\sffamily\bfseries\color{ubbblue}}
  {\Large\MakeUppercase{\chaptertitlename}~\thechapter}
  {10pt}
  {\Huge}
\titlespacing*{\chapter}{0pt}{-10pt}{28pt}

% --- Format judul section ---------------------------------------------
\titleformat{\section}
  {\normalfont\Large\bfseries\color{ubbblue}}
  {\thesection}{0.6em}{}
\titlespacing*{\section}{0pt}{2.0ex plus 0.5ex minus 0.2ex}{1.0ex}

\titleformat{\subsection}
  {\normalfont\large\bfseries\color{ubbteal}}
  {\thesubsection}{0.5em}{}
\titlespacing*{\subsection}{0pt}{1.5ex plus 0.3ex minus 0.2ex}{0.8ex}

% --- Daftar isi --------------------------------------------------------
\setcounter{tocdepth}{2}
\setlength{\cftchapnumwidth}{2.8em}
\renewcommand{\cftchapfont}{\bfseries}
\renewcommand{\cftchappagefont}{\bfseries}
\renewcommand{\cftchapleader}{\cftdotfill{\cftdotsep}}

% --- Keterangan gambar/tabel -----------------------------------------
\captionsetup{
  font=small,
  labelfont={bf,color=ubbblue},
  labelsep=period,
  skip=6pt
}

% --- Header dan footer -------------------------------------------------
\pagestyle{fancy}
\fancyhf{}
\fancyhead[L]{\small\itshape\color{textgray}Teknik Kompilasi}
\fancyhead[R]{\small\color{textgray}\thepage}
\renewcommand{\headrulewidth}{0.4pt}
\renewcommand{\footrulewidth}{0pt}

\fancypagestyle{plain}{%
  \fancyhf{}%
  \fancyfoot[C]{\small\thepage}%
  \renewcommand{\headrulewidth}{0pt}%
}

% --- Warna tautan ------------------------------------------------------
\hypersetup{
  colorlinks=true,
  linkcolor=ubbblue,
  citecolor=ubbteal,
  urlcolor=ubbteal,
  filecolor=ubbteal
}

% --- Penomoran ---------------------------------------------------------
% Bab memakai angka Romawi (Bab I, Bab II, ...).
% Section, gambar, tabel, kode, dan rumus memakai nomor bab Arab.
\renewcommand{\thechapter}{\Roman{chapter}}
\renewcommand{\thesection}{\arabic{chapter}.\arabic{section}}
\renewcommand{\thesubsection}{\thesection.\arabic{subsection}}
\renewcommand{\thefigure}{\arabic{chapter}.\arabic{figure}}
\renewcommand{\thetable}{\arabic{chapter}.\arabic{table}}
\renewcommand{\theequation}{\arabic{chapter}.\arabic{equation}}

% --- Paragraf yang lebih rapi di kolom sempit -------------------------
\setlength{\skip\footins}{10pt}

% =====================================================================
% config/commands.tex
% Perintah bersama. Jangan mendefinisikan perintah yang sama di bab.
% =====================================================================

% --- Teks dan istilah --------------------------------------------------
\newcommand{\code}[1]{\texttt{#1}}
\newcommand{\file}[1]{\texttt{#1}}
\newcommand{\eng}[1]{\textit{#1}}
\newcommand{\keyword}[1]{\textbf{#1}}

% --- Rujukan kompetensi ------------------------------------------------
\newcommand{\subcpmk}[1]{Sub-CPMK~#1}
\newcommand{\subcpmkref}[1]{Sub-CPMK~#1}
\newcommand{\cpmk}[1]{CPMK~#1}
\newcommand{\cpl}[1]{CPL~#1}

% --- Rujukan silang berbahasa Indonesia -------------------------------
\newcommand{\figref}[1]{Gambar~\ref{#1}}
\newcommand{\tabref}[1]{Tabel~\ref{#1}}
\newcommand{\lstref}[1]{Kode Program~\ref{#1}}
\newcommand{\chapref}[1]{Bab~\ref{#1}}
\newcommand{\secref}[1]{Bagian~\ref{#1}}

% --- Simbol ------------------------------------------------------------
\newcommand{\checkbox}{$\square$}
\newcommand{\stage}[1]{\textsf{\textbf{[#1]}}}

% --- Kode program ------------------------------------------------------
% Path dibungkus \subfix agar benar saat build penuh maupun mandiri.
\newcommand{\includecode}[2][]{\lstinputlisting[#1]{\subfix{#2}}}

% --- Penomoran lampiran ------------------------------------------------
\newcommand{\appendixnumbering}{%
  \renewcommand{\thechapter}{\Alph{chapter}}%
  \renewcommand{\thesection}{\Alph{chapter}.\arabic{section}}%
  \renewcommand{\thesubsection}{\thesection.\arabic{subsection}}%
  \renewcommand{\thefigure}{\Alph{chapter}.\arabic{figure}}%
  \renewcommand{\thetable}{\Alph{chapter}.\arabic{table}}%
  \renewcommand{\theequation}{\Alph{chapter}.\arabic{equation}}%
  \renewcommand{\thelstlisting}{\Alph{chapter}.\arabic{lstlisting}}%
  \renewcommand{\theHchapter}{\Alph{chapter}}%
  \renewcommand{\theHsection}{\Alph{chapter}.\arabic{section}}%
}

% =====================================================================
% config/environments.tex
% Lingkungan teorema dan kotak OBE.
% =====================================================================

% --- Teorema -----------------------------------------------------------
\theoremstyle{definition}
\newtheorem{definisi}{Definisi}[chapter]
\newtheorem{contoh}[definisi]{Contoh}

\theoremstyle{plain}
\newtheorem{teorema}[definisi]{Teorema}
\newtheorem{lemma}[definisi]{Lemma}
\newtheorem{akibat}[definisi]{Akibat}

% Nomor teorema memakai nomor bab Arab (Bab III -> Definisi 3.1).
\renewcommand{\thedefinisi}{\arabic{chapter}.\arabic{definisi}}

% --- Kotak OBE ---------------------------------------------------------
\newtcolorbox{outcomebox}[1][]{
  enhanced,
  breakable,
  colback=ubbblue!5,
  colframe=ubbblue,
  coltitle=white,
  fonttitle=\bfseries,
  title={Outcome Pembelajaran dan Sub-CPMK},
  arc=2mm,
  boxrule=0.6pt,
  left=3mm,
  right=3mm,
  top=2mm,
  bottom=2mm,
  #1
}

\newtcolorbox{infobox}[1][]{
  enhanced,
  breakable,
  colback=ubbteal!5,
  colframe=ubbteal,
  coltitle=white,
  fonttitle=\bfseries,
  title={Informasi},
  arc=2mm,
  boxrule=0.5pt,
  left=3mm,
  right=3mm,
  top=2mm,
  bottom=2mm,
  #1
}

\newtcolorbox{warningbox}[1][]{
  enhanced,
  breakable,
  colback=ubbgold!8,
  colframe=ubbgold,
  coltitle=black,
  fonttitle=\bfseries,
  title={Peringatan},
  arc=2mm,
  boxrule=0.5pt,
  left=3mm,
  right=3mm,
  top=2mm,
  bottom=2mm,
  #1
}

\newtcolorbox{tipsbox}[1][]{
  enhanced,
  breakable,
  colback=green!4,
  colframe=green!45!black,
  coltitle=white,
  fonttitle=\bfseries,
  title={Tips},
  arc=2mm,
  boxrule=0.5pt,
  left=3mm,
  right=3mm,
  top=2mm,
  bottom=2mm,
  #1
}

\newtcolorbox{reflectionbox}[1][]{
  enhanced,
  breakable,
  colback=purple!4,
  colframe=purple!55!black,
  coltitle=white,
  fonttitle=\bfseries,
  title={Refleksi},
  arc=2mm,
  boxrule=0.5pt,
  left=3mm,
  right=3mm,
  top=2mm,
  bottom=2mm,
  #1
}

\newtcolorbox{casestudy}[1][]{
  enhanced,
  breakable,
  colback=gray!4,
  colframe=gray!60!black,
  coltitle=white,
  fonttitle=\bfseries,
  title={Studi Kasus},
  arc=2mm,
  boxrule=0.5pt,
  left=3mm,
  right=3mm,
  top=2mm,
  bottom=2mm,
  #1
}

% --- Daftar cek kompetensi --------------------------------------------
\newenvironment{checklist}
  {\begin{itemize}[label=$\square$,leftmargin=*,itemsep=2pt,topsep=2pt]}
  {\end{itemize}}

% --- Daftar aktivitas --------------------------------------------------
\newenvironment{activitylist}
  {\begin{enumerate}[label=\textbf{A\arabic*.},leftmargin=*,itemsep=3pt]}
  {\end{enumerate}}

% =====================================================================
% config/listings.tex
% Gaya dan bahasa untuk kode program.
% =====================================================================

\lstdefinestyle{bookcode}{
  basicstyle=\ttfamily\small,
  keywordstyle=\color{ubbblue}\bfseries,
  commentstyle=\color{green!40!black}\itshape,
  stringstyle=\color{ubbteal},
  numbers=left,
  numberstyle=\tiny\color{textgray},
  numbersep=8pt,
  frame=single,
  rulecolor=\color{codeframe},
  backgroundcolor=\color{codebg},
  breaklines=true,
  breakatwhitespace=false,
  columns=fullflexible,
  keepspaces=true,
  showstringspaces=false,
  tabsize=2,
  captionpos=b,
  xleftmargin=1.5em,
  framexleftmargin=1.5em,
  upquote=true,
  extendedchars=false,
  inputencoding=utf8
}

\lstset{style=bookcode}

% --- Bahasa C memakai definisi bawaan listings ------------------------
% Kata kunci tambahan dapat diberikan per-listing melalui
% morekeywords={...} pada \includecode.

% --- Flex --------------------------------------------------------------
\lstdefinelanguage{Flex}{
  morekeywords={\%option,\%\{,\%\},\%\%,\%s,\%x,yylex,yylval,yytext,yyleng,yyin,yyout,REJECT,BEGIN,ECHO,YY_USER_ACTION},
  morecomment=[l]{//},
  morecomment=[s]{/*}{*/},
  morestring=[b]",
  sensitive=true
}

% --- Bison -------------------------------------------------------------
\lstdefinelanguage{Bison}{
  morekeywords={\%token,\%left,\%right,\%nonassoc,\%type,\%union,\%start,\%prec,\%expect,\%define,\%code,\%destructor,\%printer,\%parse-param,\%lex-param,\%\%},
  morecomment=[l]{//},
  morecomment=[s]{/*}{*/},
  morestring=[b]",
  sensitive=true
}

% --- EBNF --------------------------------------------------------------
\lstdefinelanguage{EBNF}{
  morekeywords={},
  morecomment=[l]{//},
  morestring=[b]",
  literate=
    {=}{{$\;=\;$}}2
    {|}{{$\;\mid\;$}}2
    {->}{{$\;\rightarrow\;$}}2
}

% --- Bash --------------------------------------------------------------
\lstdefinelanguage{bash}{
  morekeywords={if,then,else,fi,for,do,done,while,case,esac,function,echo,cd,make,gcc,flex,bison},
  morecomment=[l]{\#},
  morestring=[b]",
  morestring=[b]'
}

% --- Nama dan penomoran ------------------------------------------------
\renewcommand{\lstlistingname}{Kode Program}
\renewcommand{\lstlistlistingname}{Daftar Kode Program}

\makeatletter
\@addtoreset{lstlisting}{chapter}
\makeatother

\AtBeginDocument{%
  \renewcommand{\thelstlisting}{\arabic{chapter}.\arabic{lstlisting}}%
}

% =====================================================================
% config/tikz.tex
% Library TikZ/forest dan gaya diagram bersama.
% =====================================================================

\usetikzlibrary{
  arrows.meta,
  positioning,
  calc,
  shapes.geometric,
  shapes.misc,
  fit,
  backgrounds,
  trees,
  chains,
  automata,
  decorations.pathreplacing,
  patterns
}

\tikzset{
  phase/.style={
    rectangle, rounded corners=2pt, draw=ubbblue, fill=ubbblue!10,
    minimum width=2.6cm, minimum height=0.9cm, align=center, font=\small
  },
  data/.style={
    cylinder, shape border rotate=90, draw=ubbteal, fill=ubbteal!10,
    minimum height=1.0cm, minimum width=1.6cm, align=center, font=\small
  },
  decision/.style={
    diamond, draw=ubbgold, fill=ubbgold!15, aspect=2,
    align=center, font=\small, inner sep=1pt
  },
  arrow/.style={-{Stealth[length=2.2mm]}, thick, draw=ubbblue},
  state/.style={
    circle, draw=ubbblue, fill=ubbblue!5, minimum size=7mm, font=\small
  },
  accepting/.style={
    circle, double, double distance=1pt, draw=ubbblue,
    fill=ubbblue!5, minimum size=7mm, font=\small
  },
  memorycell/.style={
    rectangle, draw=gray, minimum width=2.2cm, minimum height=0.8cm,
    font=\small
  },
  scope/.style={
    rectangle, draw=ubbteal, dashed, rounded corners=2pt, inner sep=6pt
  }
}

\forestset{
  parsetree/.style={
    for tree={
      font=\small, s sep=3mm, l sep=7mm,
      edge={thick, draw=ubbblue}, anchor=north
    }
  },
  asttree/.style={
    for tree={
      font=\small, s sep=3mm, l sep=7mm,
      edge={thick, draw=ubbteal}, anchor=north
    }
  }
}


% =====================================================================
% metadata.tex
% Identitas buku dan metadata PDF.
% =====================================================================

\newcommand{\BookTitle}{Teknik Kompilasi}
\newcommand{\BookSubtitle}{Dari Teori Bahasa ke Implementasi Front-End Kompilator}
\newcommand{\BookEdition}{Edisi Pertama}
\newcommand{\BookYear}{2026}
\newcommand{\BookAuthor}{Cecep Suwanda, S.Si., M.Kom.}
\newcommand{\BookProgram}{Program Studi Teknik Informatika}
\newcommand{\BookFaculty}{Fakultas Teknologi Informasi}
\newcommand{\BookInstitution}{Universitas Bale Bandung}
\newcommand{\BookCode}{MWCA017}
\newcommand{\BookSemester}{Semester Ganjil}
\newcommand{\BookCourse}{Teknik Kompilasi}

\hypersetup{
  pdftitle={\BookTitle: \BookSubtitle},
  pdfauthor={\BookAuthor},
  pdfsubject={Buku ajar Teknik Kompilasi berbasis Outcome-Based Education},
  pdfkeywords={kompilator, analisis leksikal, parsing, analisis semantik, OBE},
  pdfcreator={LaTeX},
  pdfproducer={MiKTeX}
}


\begin{document}

% ---------------------------------------------------------------------
% FRONT MATTER
% ---------------------------------------------------------------------
\frontmatter
\pagestyle{plain}

% Penomoran tabel/gambar front matter tidak memakai nomor bab.
\renewcommand{\thetable}{\arabic{table}}
\renewcommand{\thefigure}{\arabic{figure}}

% =====================================================================
% frontmatter/cover.tex
% =====================================================================

\begin{titlepage}
\thispagestyle{empty}
\begin{center}

{\sffamily\bfseries\color{ubbblue}\Large BUKU AJAR}\\[5mm]

{\sffamily\bfseries\color{ubbblue}\fontsize{32}{38}\selectfont TEKNIK KOMPILASI}\\[6mm]

{\large\itshape Dari Teori Bahasa ke\\[1mm]
Implementasi Front-End Kompilator}\\[12mm]

\begin{tikzpicture}[node distance=5mm, scale=0.86, transform shape]
  \node[phase] (src) {Program\\Sumber};
  \node[phase, right=of src] (lex) {Analisis\\Leksikal};
  \node[phase, right=of lex] (syn) {Analisis\\Sintaksis};
  \node[phase, right=of syn] (sem) {Analisis\\Semantik};
  \node[phase, right=of sem] (ast) {AST dan\\Validasi};
  \draw[arrow] (src) -- (lex);
  \draw[arrow] (lex) -- (syn);
  \draw[arrow] (syn) -- (sem);
  \draw[arrow] (sem) -- (ast);
\end{tikzpicture}\\[14mm]

{\large Buku Ajar Berbasis \textit{Outcome-Based Education}}\\[18mm]

{\Large\bfseries \BookAuthor}\\[4mm]

{\normalsize
\BookProgram\\
\BookFaculty\\
\BookInstitution}\\[10mm]

{\normalsize Mata Kuliah \BookCourse\ (\BookCode)\\
\BookSemester\ --- \BookYear}

\vfill
\end{center}
\end{titlepage}

% =====================================================================
% frontmatter/copyright.tex
% =====================================================================

\thispagestyle{empty}

\vspace*{\fill}

\begin{flushleft}
{\bfseries \BookTitle: \BookSubtitle}\\[2mm]
\BookEdition\\[6mm]

\copyright\ \BookYear\ \BookAuthor\\[2mm]
\BookProgram\\
\BookFaculty\\
\BookInstitution\\[6mm]

Hak cipta dilindungi undang-undang. Dilarang memperbanyak sebagian
atau seluruh isi buku ini dalam bentuk apa pun tanpa izin tertulis
dari penulis dan penerbit.\\[6mm]

ISBN: dalam proses\\[2mm]
Katalog dalam terbitan: dalam proses\\[6mm]

Dicetak dan diterbitkan oleh\\
\BookProgram, \BookFaculty, \BookInstitution\\[2mm]
Bandung, \BookYear
\end{flushleft}

\vspace*{\fill}

\chapter*{Kata Pengantar}
\addcontentsline{toc}{chapter}{Kata Pengantar}
\markboth{Kata Pengantar}{Kata Pengantar}

Buku ajar ini disusun untuk mendukung pembelajaran mata kuliah
\textbf{Teknik Kompilasi} pada Program Studi Teknik Informatika,
Fakultas Teknologi Informasi, Universitas Bale Bandung. Mata kuliah ini
membahas bagaimana sebuah program sumber berbahasa tingkat tinggi
diterjemahkan menjadi representasi yang dapat diproses lebih lanjut
oleh mesin. Penekanan buku ini adalah \emph{front-end} kompilator:
analisis leksikal, analisis sintaksis, analisis semantik, serta
pembentukan dan validasi \emph{Abstract Syntax Tree} (AST).

Penyusunan buku mengikuti Rencana Pembelajaran Semester (RPS) berbasis
\emph{Outcome-Based Education} (OBE). Setiap bab dimulai dengan outcome
pembelajaran dan Sub-CPMK yang jelas, dilanjutkan dengan materi,
contoh, praktikum, latihan, asesmen, daftar cek kompetensi, rangkuman,
dan evaluasi. Dengan struktur ini, mahasiswa mengetahui sejak awal
kompetensi yang harus dicapai dan bukti belajar yang harus ditunjukkan.

Buku ini terdiri atas tujuh belas bab. Bab I dan Bab II memberikan
orientasi serta landasan teori. Bab III sampai Bab XVI merupakan empat
belas bab inti yang mengikuti empat belas topik RPS secara berurutan.
Bab XVII mengintegrasikan seluruh kompetensi melalui evaluasi,
refleksi, dan proyek akhir. Empat lampiran memuat rubrik penilaian,
contoh tugas dan portofolio, glosarium, serta sumber belajar tambahan.

Sepanjang semester, mahasiswa mengerjakan satu proyek berkelanjutan:
mini-kompilator untuk bahasa \textbf{Subset-C}. Proyek ini tumbuh dari
bab ke bab, mulai dari spesifikasi bahasa dan \emph{lexer}, dilanjutkan
dengan \emph{parser}, AST, tabel simbol, pemeriksaan tipe, hingga
validator dan diagnostik. Dengan cara ini, teori yang dipelajari
langsung berubah menjadi artefak perangkat lunak yang dapat diuji.

Buku ini memanfaatkan sumber-sumber terbuka yang telah tersedia,
antara lain karya Thain, Mogensen, Nystrom, Guntoro, dan catatan kuliah
Asep Juarna, serta sejumlah sumber daring yang tercantum pada daftar
pustaka. Penulis menyadari bahwa buku ajar ini masih dapat
disempurnakan. Kritik dan saran dari dosen, mahasiswa, dan pembaca
sangat diharapkan untuk perbaikan edisi berikutnya.

\begin{flushright}
Bandung, \BookYear\\
\BookAuthor
\end{flushright}

\chapter*{Ucapan Terima Kasih}
\addcontentsline{toc}{chapter}{Ucapan Terima Kasih}
\markboth{Ucapan Terima Kasih}{Ucapan Terima Kasih}

Penulis mengucapkan terima kasih kepada seluruh pihak yang telah
membantu penyusunan buku ajar ini. Ucapan terima kasih khusus
disampaikan kepada para penulis sumber terbuka yang karya-karyanya
menjadi rujukan utama: Douglas Thain, Torben \AE{}gidius Mogensen,
Robert Nystrom, Guntoro, dan Asep Juarna. Karya mereka membuat materi
teknik kompilasi dapat dipelajari secara luas dan terjangkau.

Terima kasih juga disampaikan kepada pimpinan Program Studi Teknik
Informatika, Fakultas Teknologi Informasi, Universitas Bale Bandung,
atas dukungan dalam penyusunan bahan ajar berbasis OBE. Ucapan terima
kasih kepada rekan dosen atas diskusi mengenai capaian pembelajaran,
metode asesmen, dan penyelarasan materi dengan RPS.

Penulis berterima kasih kepada mahasiswa yang telah mengikuti mata
kuliah ini. Pertanyaan, kesulitan, dan umpan balik mereka menjadi
bahan berharga untuk memperbaiki urutan materi, contoh, dan latihan.
Semoga buku ini bermanfaat bagi pembelajaran Teknik Kompilasi dan
menumbuhkan minat mahasiswa pada perancangan bahasa dan perangkat
lunak sistem.

\chapter*{Capaian Pembelajaran dan Peta Kurikulum}
\addcontentsline{toc}{chapter}{Capaian Pembelajaran dan Peta Kurikulum}
\markboth{Capaian Pembelajaran dan Peta Kurikulum}{Capaian Pembelajaran dan Peta Kurikulum}

Bagian ini merangkum capaian pembelajaran, peta Sub-CPMK, dan
keterkaitan setiap bab dengan pertemuan pada RPS. Tabel-tabel di sini
menjadi acuan untuk membaca seluruh buku dan memeriksa ketercapaian
kompetensi.

\section*{Capaian Pembelajaran Lulusan (CPL)}
\addcontentsline{toc}{section}{Capaian Pembelajaran Lulusan (CPL)}

\begin{table}[H]
\centering
\caption{CPL yang didukung mata kuliah}
\label{tab:cpl}
\begin{tabularx}{\textwidth}{@{}lX@{}}
\toprule
\textbf{Kode} & \textbf{Rumusan} \\
\midrule
CPL 02 & Memiliki pengetahuan yang memadai terkait dengan cara kerja
sistem komputer dan mampu merancang dan mengembangkan berbagai
algoritma metode untuk memecahkan masalah. \\
CPL 03 & Kemampuan merancang dan menganalisis algoritma untuk
menyelesaikan permasalahan organisasi secara optimal, serta memilih
dan menerapkannya pada bahasa pemrograman tertentu. \\
\bottomrule
\end{tabularx}
\end{table}

\section*{Capaian Pembelajaran Mata Kuliah (CPMK)}
\addcontentsline{toc}{section}{Capaian Pembelajaran Mata Kuliah (CPMK)}

\begin{table}[H]
\centering
\caption{CPMK dan fokus kompetensi}
\label{tab:cpmk}
\begin{tabularx}{\textwidth}{@{}lX@{}}
\toprule
\textbf{Kode} & \textbf{Rumusan} \\
\midrule
CPMK 01 & Menjelaskan tahapan kompilasi (lexical, syntax, semantic
analysis). \\
CPMK 02 & Memahami konsep Otomata Hingga (Finite Automata). \\
CPMK 03 & Merancang tata bahasa (grammar) bebas konteks untuk bahasa
sederhana. \\
CPMK 04 & Mengimplementasikan parser sederhana untuk validasi kode. \\
\bottomrule
\end{tabularx}
\end{table}

\section*{Peta Sub-CPMK}
\addcontentsline{toc}{section}{Peta Sub-CPMK}

{\small
\begin{longtable}{@{}p{1.2cm}p{7.8cm}p{1.6cm}p{3.5cm}@{}}
\caption{Peta sembilan belas Sub-CPMK ke bab dan asesmen}
\label{tab:subcpmk}\\
\toprule
\textbf{Kode} & \textbf{Rumusan} & \textbf{Bab} & \textbf{Asesmen} \\
\midrule
\endfirsthead
\toprule
\textbf{Kode} & \textbf{Rumusan} & \textbf{Bab} & \textbf{Asesmen} \\
\midrule
\endhead
\bottomrule
\endfoot
1.1 & Menjelaskan perbedaan interpreter dan compiler serta peran
assembler, interpreter, dan compiler sebagai translator, dengan
menyebut contoh masukan dan keluaran masing-masing. (CPL 02) & III &
Kuis Konsep, UAS Tulis \\
1.2 & Menjelaskan fase-fase kompilator (analisis leksikal, sintaksis,
semantik, dan pembangkitan kode antara) beserta peran tabel simbol
dan penanganan kesalahan. (CPL 02) & III & Kuis Konsep, UAS Tulis \\
1.3 & Membuat tabel penelusuran fase kompilasi untuk satu program C
sederhana, yang memuat keluaran antara berupa token, pohon sintaks,
dan kode antara. (CPL 02) & III & UTS Praktik \\
1.4 & Membuat tabel pemetaan hierarki Chomsky (Tipe 3--0) ke jenis
automata dan menguji klasifikasi lima contoh grammar ke kelas yang
tepat. (CPL 02) & IV & Kuis Konsep, UTS Praktik \\
1.5 & Membuat tabel aturan semantik berbasis \textit{syntax-directed
definition} untuk ekspresi dan pernyataan subset bahasa C, lalu
menguji keluarannya pada pohon sintaks. (CPL 02) & XIV & Praktikum 4 \\
1.6 & Mengimplementasikan \textit{symbol table} bersarang dan pemeriksa
tipe dalam bahasa C, dengan berkas kode dan keluaran diagnosis sebagai
bukti. (CPL 02) & XV & Praktikum 4 \\
1.7 & Menguji dan men-\textit{debug} penanganan kesalahan semantik
sehingga program menghasilkan pesan kesalahan yang menyebut baris,
nama, dan jenis pelanggaran. (CPL 02) & XVI & Praktikum 4, Proyek
Akhir \\
2.1 & Membuat tabel ekspresi reguler untuk setiap kategori token subset
bahasa C (kata kunci, \textit{identifier}, bilangan, operator, dan
pemisah). (CPL 02) & V & Praktikum 1 \\
2.2 & Mengimplementasikan simulator NFA dan DFA dalam bahasa C untuk
mengenali token, dengan berkas kode dan keluaran pengujian sebagai
bukti. (CPL 02) & VI & Praktikum 1, UTS Praktik \\
2.3 & Membuat tabel konversi NFA ke DFA dan tabel minimalisasi DFA,
serta menguji bahwa keduanya menerima bahasa token yang sama.
(CPL 02) & VI & Praktikum 1, UTS Praktik \\
2.4 & Mengimplementasikan \textit{lexer} subset bahasa C dalam dua
versi, yaitu \textit{hand-written} dan Flex, lalu menguji keluarannya
pada berkas masukan uji. (CPL 02) & VII & Praktikum 1, Proyek Akhir \\
3.1 & Membuat berkas grammar BNF dan EBNF untuk subset bahasa C, serta
menguji derivasi dan pohon sintaks yang dihasilkan. (CPL 03) & VIII &
Praktikum 2, UTS Praktik \\
3.2 & Menguji grammar subset bahasa C dan men-\textit{debug}
ambiguitas, \textit{left recursion}, serta \textit{left factoring}
sampai berkas grammar bebas konflik. (CPL 03) & IX & Praktikum 2 \\
3.3 & Membuat tabel himpunan FIRST dan FOLLOW untuk grammar subset
bahasa C dan menguji keputusan parser LL(1) berdasarkan tabel
tersebut. (CPL 03) & X & Praktikum 2, UTS Praktik \\
3.4 & Membuat berkas grammar bebas konteks untuk subset bahasa C yang
mencakup deklarasi, ekspresi, dan pernyataan, sebagai dasar parser dan
proyek akhir. (CPL 03) & VIII & Praktikum 2, Proyek Akhir \\
4.1 & Mengimplementasikan \textit{recursive descent parser} dan
\textit{predictive parser} LL(1) dalam bahasa C, dengan berkas kode
dan keluaran validasi sebagai bukti. (CPL 03) & X & Praktikum 2 \\
4.2 & Membuat tabel aksi dan \textit{goto} untuk parser SLR/LALR, lalu
menguji dan men-\textit{debug} konflik \textit{shift-reduce} atau
\textit{reduce-reduce} yang muncul. (CPL 03) & XI--XII & Praktikum 3,
UTS Praktik \\
4.3 & Mengimplementasikan parser subset bahasa C dengan Flex dan Bison
beserta aksi semantik, dengan berkas kode dan keluaran program sebagai
bukti. (CPL 03) & XIII & Praktikum 3, Proyek Akhir \\
4.4 & Membuat AST dan validator kode untuk subset bahasa C, lalu
menguji keluaran pesan kesalahan sintaksis yang informatif.
(CPL 03) & XIII, XVI & Praktikum 3, Proyek Akhir \\
\end{longtable}
}

\section*{Peta Pertemuan dengan Bab}
\addcontentsline{toc}{section}{Peta Pertemuan dengan Bab}

{\small
\begin{longtable}{@{}p{0.9cm}p{5.8cm}p{1.6cm}p{2.6cm}p{2.7cm}@{}}
\caption{Peta enam belas pertemuan RPS ke bab}
\label{tab:peta-pertemuan}\\
\toprule
\textbf{Pert.} & \textbf{Topik} & \textbf{Bab} & \textbf{Sub-CPMK} &
\textbf{Asesmen} \\
\midrule
\endfirsthead
\toprule
\textbf{Pert.} & \textbf{Topik} & \textbf{Bab} & \textbf{Sub-CPMK} &
\textbf{Asesmen} \\
\midrule
\endhead
\bottomrule
\endfoot
1 & Pengenalan kompilasi dan arsitektur kompilator & III & 1.1--1.3 &
Kuis Konsep \\
2 & Teori bahasa formal dan hierarki Chomsky & IV & 1.4 &
Kuis Konsep \\
3 & Analisis leksikal: token, lexeme, pattern, regex & V & 2.1 &
Praktikum 1 \\
4 & Finite automata: NFA, DFA, konversi, minimalisasi & VI & 2.2--2.3 &
Praktikum 1 \\
5 & Implementasi lexical analyzer & VII & 2.4 & Praktikum 1 \\
6 & Context-free grammar: BNF, EBNF, derivasi, parse tree & VIII &
3.1, 3.4 & Praktikum 2 \\
7 & Ambiguitas, left recursion, left factoring & IX & 3.2 &
Praktikum 2 \\
8 & Ujian Tengah Semester (UTS) & --- & CPMK 01--03 & UTS Praktik \\
9 & Top-down parsing: recursive descent, LL(1), FIRST/FOLLOW & X &
3.3, 4.1 & Praktikum 2 \\
10 & Bottom-up parsing: shift-reduce, operator precedence & XI & 4.2 &
Praktikum 3 \\
11 & LR parsing: canonical LR, SLR, LALR, konflik & XII & 4.2 &
Praktikum 3 \\
12 & Parser generator dan semantic actions & XIII & 4.3--4.4 &
Praktikum 3 \\
13 & Analisis semantik: tabel aturan semantik (SDD) & XIV & 1.5 &
Praktikum 4 \\
14 & Symbol table, scope management, type checking & XV & 1.6 &
Praktikum 4 \\
15 & AST, validasi kode, penanganan error & XVI & 4.4, 1.7 &
Praktikum 4, Proyek Akhir \\
16 & Ujian Akhir Semester (UAS) & --- & CPMK 01--04 &
UAS Tulis, Proyek Akhir \\
\end{longtable}
}

\section*{Bobot Penilaian}
\addcontentsline{toc}{section}{Bobot Penilaian}

\begin{table}[H]
\centering
\caption{Komponen penilaian dan bobot}
\label{tab:bobot-penilaian}
\begin{tabularx}{\textwidth}{@{}lXr@{}}
\toprule
\textbf{Komponen} & \textbf{Cakupan} & \textbf{Bobot} \\
\midrule
Praktikum 1 & Lexer (Sub-CPMK 2.1--2.4) & 12\% \\
Praktikum 2 & Grammar dan parser (Sub-CPMK 3.1--3.4, 4.1) & 12\% \\
Praktikum 3 & Tabel SLR/LALR, Bison, dan AST (Sub-CPMK 4.2--4.4) & 12\% \\
Praktikum 4 & Symbol table dan analisis semantik (Sub-CPMK 1.5--1.7) & 12\% \\
Proyek Akhir & Front-end mini-compiler subset C & 30\% \\
UTS Praktik & Ujian praktik (Sub-CPMK 1.3, 2.2, 2.3, 3.1, 3.3, 4.2) & 14\% \\
Kuis Konsep & Konsep kompilasi (Sub-CPMK 1.1, 1.2, 1.4) & 3\% \\
UAS Tulis & Konsep dan analisis (Sub-CPMK 1.1, 1.2, 1.4) & 5\% \\
\midrule
\textbf{Total} & & \textbf{100\%} \\
\bottomrule
\end{tabularx}
\end{table}

Komposisi bobot adalah 92\% praktik (Praktikum 1--4, Proyek Akhir, dan
UTS Praktik) dan 8\% teori (Kuis Konsep dan UAS Tulis).

\chapter*{Cara Menggunakan Buku Ini}
\addcontentsline{toc}{chapter}{Cara Menggunakan Buku Ini}
\markboth{Cara Menggunakan Buku Ini}{Cara Menggunakan Buku Ini}

\section*{Anatomi Bab}
\addcontentsline{toc}{section}{Anatomi Bab}

Setiap bab inti memiliki urutan yang sama:
\begin{enumerate}[label=\arabic*.,leftmargin=*]
  \item \textbf{Outcome Pembelajaran dan Sub-CPMK}: kompetensi yang
        harus dicapai dan kode Sub-CPMK yang terkait.
  \item \textbf{Materi Pokok}: peta isi bab.
  \item \textbf{Lima bagian materi}: uraian konsep, definisi, contoh,
        dan ilustrasi.
  \item \textbf{Contoh}: soal berpenyelesaian lengkap.
  \item \textbf{Aktivitas Pembelajaran dan Praktikum}: latihan
        terbimbing dan increment proyek.
  \item \textbf{Latihan dan Refleksi}: soal latihan bertingkat dan
        pertanyaan refleksi.
  \item \textbf{Asesmen}: instrumen dan rubrik yang mengukur Sub-CPMK.
  \item \textbf{Checklist Pencapaian Kompetensi}: bukti yang harus
        ditunjukkan mahasiswa.
  \item \textbf{Rangkuman}: ringkasan dan jembatan ke bab berikutnya.
  \item \textbf{Evaluasi}: soal evaluasi akhir bab dan rubrik singkat.
\end{enumerate}

\section*{Arti Kotak}
\addcontentsline{toc}{section}{Arti Kotak}

\begin{infobox}
Kotak \textbf{Informasi} memuat konsep pendukung, konteks sejarah,
atau keterkaitan dengan bidang lain.
\end{infobox}

\begin{warningbox}
Kotak \textbf{Peringatan} memuat kesalahan umum, jebakan konseptual,
atau batasan yang perlu diperhatikan.
\end{warningbox}

\begin{tipsbox}
Kotak \textbf{Tips} memuat strategi praktis untuk belajar, menulis
kode, atau mengerjakan ujian.
\end{tipsbox}

\begin{reflectionbox}
Kotak \textbf{Refleksi} mengajak mahasiswa menilai pemahaman dan
kebiasaan belajarnya sendiri.
\end{reflectionbox}

\section*{Proyek Mini-Kompilator}
\addcontentsline{toc}{section}{Proyek Mini-Kompilator}

Proyek berkelanjutan pada buku ini adalah mini-kompilator untuk
\textbf{Subset-C}. Spesifikasi bahasa, kode sumber, contoh program,
dan tes tersedia di direktori \file{code/subsetc}. Proyek dibagi
menjadi enam milestone (M0--M5) yang dikerjakan pada bab-bab tertentu.
Setiap bab praktikum membuka dengan tabel artefak: berkas, status awal,
pekerjaan yang harus dilakukan, dan definisi selesai.

Mahasiswa disarankan menyimpan proyek dalam repositori versi dan
melakukan commit pada setiap milestone. Kode yang ditampilkan di buku
adalah kutipan dari berkas nyata; nomor baris pada buku sesuai dengan
nomor baris berkas.

\section*{Dua Jalur Belajar}
\addcontentsline{toc}{section}{Dua Jalur Belajar}

\textbf{Jalur teori} menekankan definisi, pembuktian, dan analisis.
Mahasiswa membaca bab secara berurutan, mengerjakan latihan teori, dan
menggunakan praktikum sebagai ilustrasi. Jalur ini cocok bagi mahasiswa
yang ingin memperdalam dasar formal.

\textbf{Jalur praktik} menekankan implementasi. Mahasiswa membaca bab
secukupnya, langsung mengerjakan praktikum, lalu kembali ke teori ketika
menemui kesulitan. Jalur ini cocok bagi mahasiswa yang belajar paling
baik melalui kode dan eksperimen.

Kedua jalur bertemu pada proyek akhir dan UAS. Tidak ada jalur yang
lebih mudah: keduanya menuntut pemahaman konsep dan kemampuan
implementasi.

\section*{Konvensi Notasi}
\addcontentsline{toc}{section}{Konvensi Notasi}

\begin{itemize}[leftmargin=*]
  \item Simbol nonterminal ditulis miring, misalnya \(E\), \(T\),
        dan \(F\).
  \item Simbol terminal ditulis dengan huruf tebal atau dalam tanda
        kutip, misalnya \code{id}, \code{+}, dan \code{if}.
  \item String kosong dilambangkan dengan \(\varepsilon\).
  \item Turunan ditulis dengan \(\Rightarrow\), dan turunan nol atau
        lebih dengan \(\Rightarrow^{*}\).
  \item Nama berkas dan perintah ditulis dengan huruf mesin, misalnya
        \file{lexer.c} dan \code{pdflatex}.
  \item Istilah Inggris dimiringkan pada kemunculan pertama, misalnya
        \eng{parser}, lalu bentuk Indonesia dipakai konsisten.
\end{itemize}

\section*{Mengompilasi Bab Secara Mandiri}
\addcontentsline{toc}{section}{Mengompilasi Bab Secara Mandiri}

Setiap bab dapat dikompilasi secara mandiri. Dari direktori bab,
jalankan:
\begin{lstlisting}[language=bash]
pdflatex chapter.tex
pdflatex chapter.tex
\end{lstlisting}
Rujukan ke bab lain yang belum ditulis akan tampil sebagai \code{??}.
Ini normal dan tidak mengganggu pemeriksaan tata letak. Untuk
mengompilasi seluruh buku, gunakan \file{scripts/build-all.bat}.


\tableofcontents
\cleardoublepage
\listoffigures
\cleardoublepage
\listoftables
\cleardoublepage
\lstlistoflistings
\cleardoublepage

% ---------------------------------------------------------------------
% MAIN MATTER
% Bab ditambahkan bertahap sesuai fase penulisan.
% ---------------------------------------------------------------------
\mainmatter
\pagestyle{fancy}

% Kembalikan penomoran berbasis bab untuk isi utama.
\renewcommand{\thetable}{\arabic{chapter}.\arabic{table}}
\renewcommand{\thefigure}{\arabic{chapter}.\arabic{figure}}

\subfile{chapters/01/chapter}
\subfile{chapters/02/chapter}
\subfile{chapters/03/chapter}
\subfile{chapters/04/chapter}
\subfile{chapters/05/chapter}
\subfile{chapters/06/chapter}
\subfile{chapters/07/chapter}
\subfile{chapters/08/chapter}
\subfile{chapters/09/chapter}
\subfile{chapters/10/chapter}
\subfile{chapters/11/chapter}
\subfile{chapters/12/chapter}
\subfile{chapters/13/chapter}
\subfile{chapters/14/chapter}
\subfile{chapters/15/chapter}
\subfile{chapters/16/chapter}
\subfile{chapters/17/chapter}

% ---------------------------------------------------------------------
% APPENDICES
% ---------------------------------------------------------------------
\appendix
\appendixnumbering

\chapter{Rubrik Penilaian Lengkap}
\label{app:rubrik}

Lampiran ini memuat rubrik penilaian yang dipakai sepanjang semester.
Bobot total seluruh komponen adalah 100\% dan mengikuti RPS.

\section{Bobot Komponen Penilaian}
\label{app:rubrik-bobot}

\begin{table}[H]
\centering
\caption{Komponen penilaian dan bobotnya}
\label{tab:app-bobot}
\begin{tabularx}{\textwidth}{@{}lXr@{}}
\toprule
\textbf{Komponen} & \textbf{Cakupan} & \textbf{Bobot} \\
\midrule
Kuis & Sub-CPMK 1.1--1.3 & 5\% \\
Tugas Teori & Grammar, automata, parsing & 8\% \\
Praktikum 1 & Lexer & 8\% \\
Praktikum 2 & Parser & 8\% \\
Praktikum 3 & Parser generator dan AST & 7\% \\
Praktikum 4 & Symbol table dan analisis semantik & 8\% \\
UTS & CPMK 01--03 & 20\% \\
Project Final & Mini compiler front-end & 20\% \\
UAS & CPMK 01--04 & 16\% \\
\midrule
\textbf{Total} & & \textbf{100\%} \\
\bottomrule
\end{tabularx}
\end{table}

\section{Rubrik Praktikum}
\label{app:rubrik-praktikum}

\begin{table}[H]
\centering
\caption{Rubrik umum Praktikum 1--4}
\label{tab:app-rubrik-praktikum}
\begin{tabularx}{\textwidth}{@{}lXr@{}}
\toprule
\textbf{Aspek} & \textbf{Indikator} & \textbf{Bobot} \\
\midrule
Kebenaran & Program memenuhi spesifikasi dan lulus tes & 40\% \\
Kualitas kode & Keterbacaan, struktur, penamaan, komentar & 30\% \\
Kelengkapan & Semua fitur wajib diimplementasikan & 20\% \\
Pengujian & Tes mencakup kasus normal dan kasus batas & 10\% \\
\bottomrule
\end{tabularx}
\end{table}

\section{Konversi Nilai}
\label{app:rubrik-nilai}

\begin{table}[H]
\centering
\caption{Konversi nilai akhir}
\label{tab:app-nilai}
\begin{tabular}{@{}ll@{}}
\toprule
\textbf{Rentang} & \textbf{Nilai} \\
\midrule
85--100 & A \\
70--84 & B \\
60--69 & C \\
50--59 & D \\
$<$50 & E \\
\bottomrule
\end{tabular}
\end{table}

\chapter{Contoh Tugas, Portofolio, dan Rubrik Refleksi}
\label{app:tugas}

Lampiran ini memuat contoh lembar tugas, template laporan, dan rubrik
refleksi yang dapat langsung dipakai atau disesuaikan.

\section{Lembar Tugas Praktikum}
\label{app:tugas-praktikum}

\begin{table}[H]
\centering
\caption{Format lembar tugas praktikum}
\label{tab:app-tugas-format}
\begin{tabularx}{\textwidth}{@{}lX@{}}
\toprule
\textbf{Bagian} & \textbf{Isi} \\
\midrule
Identitas & Nama, NIM, kelas, tanggal \\
Tujuan & Sub-CPMK yang diukur \\
Deskripsi & Konteks dan masalah yang diselesaikan \\
Keluaran & Berkas, laporan, demo \\
Kriteria & Rubrik dan bobot \\
Batas waktu & Tanggal dan waktu pengumpulan \\
Integritas & Pernyataan orisinalitas \\
\bottomrule
\end{tabularx}
\end{table}

\section{Template Laporan Proyek}
\label{app:tugas-laporan}

Laporan proyek memuat:
\begin{enumerate}[label=\arabic*.,leftmargin=*]
  \item \textbf{Pendahuluan}: tujuan, ruang lingkup, dan batasan.
  \item \textbf{Spesifikasi bahasa}: token, grammar, dan aturan
        semantik.
  \item \textbf{Arsitektur}: modul, alur data, dan keputusan desain.
  \item \textbf{Implementasi}: lexer, parser, AST, type checker,
        validator.
  \item \textbf{Pengujian}: kasus uji, hasil, dan cakupan.
  \item \textbf{Keterbatasan}: fitur yang belum didukung.
  \item \textbf{Refleksi}: pembelajaran dan rencana perbaikan.
  \item \textbf{Lampiran}: kode, keluaran, dan tautan repositori.
\end{enumerate}

\section{Template Jurnal Refleksi}
\label{app:tugas-jurnal}

\begin{table}[H]
\centering
\caption{Template jurnal mingguan}
\label{tab:app-jurnal}
\begin{tabularx}{\textwidth}{@{}lX@{}}
\toprule
\textbf{Pertanyaan} & \textbf{Jawaban} \\
\midrule
Apa yang dipelajari? & \\
Apa yang berhasil? & \\
Apa yang sulit? & \\
Apa yang akan dilakukan? & \\
Bukti/tautan & \\
\bottomrule
\end{tabularx}
\end{table}

\section{Rubrik Refleksi}
\label{app:tugas-refleksi}

\begin{table}[H]
\centering
\caption{Rubrik penilaian refleksi}
\label{tab:app-rubrik-refleksi}
\begin{tabularx}{\textwidth}{@{}lXr@{}}
\toprule
\textbf{Aspek} & \textbf{Indikator} & \textbf{Bobot} \\
\midrule
Kedalaman & Menjelaskan penyebab, bukan hanya kejadian & 30\% \\
Keterkaitan & Menghubungkan pengalaman dengan konsep & 25\% \\
Kejujuran & Mengakui kesulitan dan kesalahan & 20\% \\
Rencana & Tindakan konkret dan terukur & 25\% \\
\bottomrule
\end{tabularx}
\end{table}

\section{Contoh Entri Portofolio}
\label{app:tugas-portofolio}

\begin{casestudy}
\textbf{Judul:} Menangani dangling else pada parser Subset-C.

\textbf{Masalah:} Parser awal memasangkan \code{else} ke \code{if}
luar, sehingga program \code{if (a) if (b) x; else y;} berperilaku
salah.

\textbf{Tindakan:} Mengubah grammar dan menambahkan uji kasus.
Parser memilih shift pada token \code{else}, sesuai aturan else
terdekat.

\textbf{Bukti:} Berkas uji \file{if\_else.sc} dan keluaran AST.

\textbf{Refleksi:} Ambiguitas tidak cukup diselesaikan dengan
intuisi; perlu aturan eksplisit dan pengujian.
\end{casestudy}

\chapter{Glosarium}
\label{app:glosarium}

{\small
\begin{longtable}{@{}p{3.8cm}p{10.8cm}@{}}
\caption{Istilah penting dalam buku ini}
\label{tab:app-glosarium}\\
\toprule
\textbf{Istilah} & \textbf{Definisi} \\
\midrule
\endfirsthead
\toprule
\textbf{Istilah} & \textbf{Definisi} \\
\midrule
\endhead
\bottomrule
\endfoot
Alfabet & Himpunan berhingga simbol yang membentuk string. \\
Ambiguitas & Keadaan ketika satu string memiliki lebih dari satu
pohon urai. \\
Assembler & Translator dari bahasa assembly ke kode mesin. \\
Associativity & Aturan pengelompokan operator dengan precedence
sama. \\
AST & \emph{Abstract Syntax Tree}; representasi pohon yang
menghilangkan detail sintaksis. \\
Atribut & Informasi tambahan pada simpul tata bahasa, misalnya tipe
atau nilai. \\
Attribute grammar & Tata bahasa dengan atribut dan aturan semantik. \\
Back end & Bagian kompilator yang menghasilkan kode target. \\
BNF & \emph{Backus-Naur Form}; notasi tata bahasa bebas konteks. \\
Bootstrap & Proses membangun compiler menggunakan bahasa yang
dikompilasi. \\
CFG & \emph{Context-Free Grammar}; tata bahasa bebas konteks. \\
Chomsky & Hierarki kelas tata bahasa dan bahasa. \\
Closure & Himpunan item yang ditambahkan secara transitif. \\
Coercion & Konversi tipe implisit. \\
Compiler & Translator dari bahasa tingkat tinggi ke bahasa target. \\
Cross-compiler & Compiler yang menghasilkan kode untuk platform
berbeda. \\
Dangling else & Ambiguitas pemasangan \code{else} dengan \code{if}. \\
Derivasi & Urutan penerapan aturan produksi. \\
DFA & \emph{Deterministic Finite Automaton}. \\
Diagnostik & Pesan kesalahan atau peringatan dengan lokasi. \\
EBNF & \emph{Extended BNF}; BNF dengan operator pengulangan dan
opsional. \\
Epsilon closure & Himpunan state yang dicapai melalui transisi
epsilon. \\
Error recovery & Strategi melanjutkan analisis setelah kesalahan. \\
FIRST & Himpunan terminal yang dapat muncul pertama dari suatu
simbol. \\
FOLLOW & Himpunan terminal yang dapat muncul setelah suatu
nonterminal. \\
Front end & Bagian kompilator yang menganalisis program sumber. \\
Grammar & Empat-tupel nonterminal, terminal, aturan, dan simbol
awal. \\
Handle & Substring yang dapat direduksi pada shift-reduce. \\
Inherited & Atribut yang dihitung dari orang tua atau saudara. \\
Interpreter & Program yang mengeksekusi program sumber secara
langsung. \\
IR & \emph{Intermediate Representation}; representasi antara. \\
Item & Aturan produksi dengan titik posisi. \\
Kata kunci & Identifier yang dicadangkan oleh bahasa. \\
Kleene star & Operasi pengulangan nol atau lebih. \\
LALR & \emph{Look-Ahead LR}; LR(1) dengan penggabungan state. \\
Left factoring & Mengeluarkan prefiks umum dari alternatif aturan. \\
Left recursion & Aturan yang memulai ruas kanan dengan nonterminal
kiri. \\
Lexeme & String konkret yang cocok dengan token. \\
Lexer & Fase yang mengubah karakter menjadi token. \\
LL(1) & Parser top-down dengan satu token lookahead. \\
LR & Parser bottom-up dengan lookahead. \\
L-value & Ekspresi yang dapat menjadi target assignment. \\
Minimalisasi & Menggabungkan state DFA yang ekuivalen. \\
NFA & \emph{Nondeterministic Finite Automaton}. \\
Nonterminal & Simbol yang dapat diturunkan lebih lanjut. \\
One-pass & Kompilator yang membaca masukan satu kali. \\
Operator precedence & Parser bottom-up berbasis relasi precedence. \\
Parse tree & Pohon yang merepresentasikan derivasi lengkap. \\
Parser & Fase yang memeriksa struktur sintaksis. \\
Parser generator & Alat yang menghasilkan parser dari grammar. \\
Pattern & Aturan yang mendefinisikan bentuk lexeme. \\
PDA & \emph{Pushdown Automaton}; otomata dengan stack. \\
Precedence & Tingkat pengikatan operator. \\
Predictive parser & Parser top-down tanpa backtracking. \\
Produksi & Aturan penulisan ulang pada tata bahasa. \\
Recursive descent & Parser top-down dengan satu fungsi per
nonterminal. \\
Reduce & Mengganti handle dengan nonterminal. \\
Regular expression & Notasi untuk bahasa regular. \\
Scope & Wilayah keberlakuan deklarasi. \\
SDD & \emph{Syntax-Directed Definition}. \\
Semantic action & Kode yang dijalankan saat aturan direduksi. \\
Sentential form & String campuran terminal dan nonterminal dalam
derivasi. \\
Shadowing & Identifier dalam menutupi identifier luar. \\
Shift & Mendorong token ke stack. \\
SLR & \emph{Simple LR}. \\
State & Konfigurasi mesin abstrak. \\
Subset construction & Konversi NFA ke DFA. \\
Symbol table & Struktur yang memetakan nama ke informasi simbol. \\
Synthesized & Atribut yang dihitung dari anak. \\
Tata bahasa & Aturan pembentuk bahasa. \\
Terminal & Simbol yang tidak dapat diturunkan lebih lanjut. \\
Thompson & Konstruksi NFA dari ekspresi reguler. \\
Token & Kategori leksikal hasil lexer. \\
Translation scheme & Tata bahasa dengan aksi semantik tersisip. \\
Type checking & Pemeriksaan kompatibilitas tipe. \\
Viable prefix & Prefiks stack yang masih dapat menjadi bentuk
kalimat. \\
Visitor & Pola traversal yang memisahkan operasi dari struktur. \\
\end{longtable}
}

\chapter{Referensi Tambahan dan Sumber Belajar}
\label{app:referensi}

\section{Referensi RPS}
\label{app:referensi-rps}

\begin{itemize}[leftmargin=*]
  \item \textbf{[1] Thain (2024)}: pengantar compiler modern dengan
        proyek B-Minor; sangat baik untuk scanning, parsing, dan AST.
  \item \textbf{[2] Mogensen (2010)}: teori ringkas dan padat;
        kuat untuk automata, parsing, dan type checking.
  \item \textbf{[3] Bergmann (2024)}: teori, alat, dan contoh;
        berguna untuk latihan terstruktur.
  \item \textbf{[4] Bornat (2008)}: panduan praktis menulis
        compiler; berguna untuk perspektif implementasi.
  \item \textbf{[5] Critchlow dan Eck (2024)}: dasar teori
        komputasi dan automata.
  \item \textbf{[6] Srikant}: catatan NPTEL; ringkas dan berorientasi
        ujian.
  \item \textbf{[7] MIT OCW 6.035}: kuliah dan tugas compiler.
  \item \textbf{[8] Stanford CS143}: slide dan handout parsing.
  \item \textbf{[9] Saarland}: struktur compiler.
  \item \textbf{[10] McGill CS520}: handout top-down parsing.
  \item \textbf{[11] Edinburgh}: compiling techniques.
  \item \textbf{[12] Flex Manual}: referensi lexer generator.
  \item \textbf{[13] GNU Bison Manual}: referensi parser generator.
  \item \textbf{[14] ANTLR 4}: dokumentasi parser generator.
  \item \textbf{[15] Wikibooks Compiler Construction}: ringkasan
        praktis.
  \item \textbf{[16] Wikibooks Formal Languages}: teori bahasa dan
        automata.
\end{itemize}

\section{Karya Lokal}
\label{app:referensi-lokal}

\begin{itemize}[leftmargin=*]
  \item Nystrom (2021), \emph{Crafting Interpreters}: implementasi
        interpreter dan VM, sangat baik untuk scanning dan parsing.
  \item Guntoro (2024), \emph{Teknik Kompilasi}: istilah Indonesia,
        penanganan kesalahan, dan tabel informasi.
  \item Juarna, \emph{Catatan Kuliah Teknik Kompilasi}: ringkasan
        leksikal, sintaksis, semantik, dan kode antara.
\end{itemize}

\section{Alat dan Sumber Daring}
\label{app:referensi-alat}

\begin{itemize}[leftmargin=*]
  \item Compiler Explorer (\code{godbolt.org}): melihat kode assembly
        dari berbagai compiler.
  \item regex101 (\code{regex101.com}): menguji ekspresi reguler.
  \item Flex manual (\code{westes.github.io/flex/manual/}).
  \item GNU Bison manual (\code{gnu.org/software/bison/manual/}).
  \item ANTLR Lab (\code{lab.antlr.org}).
  \item LLVM Kaleidoscope (\code{llvm.org/docs/tutorial/}).
  \item TinyCC (\code{bellard.org/tcc/}).
\end{itemize}

\section{Jalur Bacaan}
\label{app:referensi-jalur}

\begin{table}[H]
\centering
\caption{Jalur bacaan menurut minat}
\label{tab:app-jalur}
\begin{tabularx}{\textwidth}{@{}l X@{}}
\toprule
\textbf{Minat} & \textbf{Bacaan} \\
\midrule
Dasar & Thain Bab 1--5; Mogensen Bab 1--3 \\
Lexer & Thain Bab 3; Flex manual \\
Parser & Thain Bab 4--5; Bison manual; Stanford CS143 \\
Semantik & Thain Bab 7; Mogensen Bab 6 \\
Praktik & Nystrom; Juarna; Guntoro \\
Teori & Critchlow dan Eck; Wikibooks Formal Languages \\
\bottomrule
\end{tabularx}
\end{table}

\section{Peta Sitasi}
\label{app:referensi-peta}

\begin{table}[H]
\centering
\caption{Kunci bibliografi dan nomor RPS}
\label{tab:app-sitasi}
\begin{tabular}{@{}ll@{}}
\toprule
\textbf{Kunci} & \textbf{Nomor RPS} \\
\midrule
\code{thain2024} & [1] \\
\code{mogensen2010} & [2] \\
\code{bergmann2024} & [3] \\
\code{bornat2008} & [4] \\
\code{critchlow2024} & [5] \\
\code{srikant} & [6] \\
\code{mitocw2010} & [7] \\
\code{stanfordcs143} & [8] \\
\code{saarland2013} & [9] \\
\code{mcgill2026} & [10] \\
\code{edinburgh2024} & [11] \\
\code{flexmanual} & [12] \\
\code{bisonmanual} & [13] \\
\code{antlr4} & [14] \\
\code{wikibookscc} & [15] \\
\code{wikibookstfl} & [16] \\
\bottomrule
\end{tabular}
\end{table}


% ---------------------------------------------------------------------
% BACK MATTER
% ---------------------------------------------------------------------
\backmatter
\nocite{*}
\renewcommand{\bibname}{Daftar Pustaka}
\bibliographystyle{apalike}
\bibliography{references}
\chapter*{Tentang Penulis}
\addcontentsline{toc}{chapter}{Tentang Penulis}
\markboth{Tentang Penulis}{Tentang Penulis}

\textbf{Cecep Suwanda, S.Si., M.Kom.} adalah dosen pada Program Studi
Teknik Informatika, Fakultas Teknologi Informasi, Universitas Bale
Bandung. Ia mengampu mata kuliah Teknik Kompilasi serta mata kuliah
lain di bidang rekayasa perangkat lunak dan sistem komputasi.

Penulis menyusun buku ajar ini untuk mendukung pembelajaran berbasis
\textit{Outcome-Based Education} (OBE) dan mendekatkan teori bahasa
formal dengan praktik pembangunan front-end kompilator. Buku ini
dikembangkan dari pengalaman mengajar, diskusi dengan mahasiswa, dan
kajian sumber-sumber terbuka yang tercantum pada daftar pustaka.

Saran, koreksi, dan masukan untuk perbaikan edisi berikutnya dapat
disampaikan melalui Program Studi Teknik Informatika, Fakultas
Teknologi Informasi, Universitas Bale Bandung.


\end{document}
